%% ma: L10-B15-0088
%% lop: 10
%% chuong: 6
%% bai: 15
%% dang: 10.15.6
%% loai: TLN
%% muc: TH
%% dap_an: "3"
%% nguon: "Ảnh giáo viên gửi 25/11/2026: Bài tập trắc nghiệm – Toán 10 Bài 15 Hàm số (3 phần), Phần III Câu 5"
%% kien_thuc: "Lớp 10 Bài 15 (giá trị của hàm số); phương trình bậc hai (lớp 9)"
%% trang_thai: chinh-thuc
%% ghi_chu: "Đơn vị đáp số: giây."
\begin{ex}
Một quả bóng được ném vào không trung có chiều cao tính từ lúc bắt đầu ném ra được cho bởi công thức $h(t)=-t^2+2t+3$ (tính bằng mét), $t$ là thời gian tính bằng giây ($t\ge0$). Hãy tính xem sau bao lâu quả bóng sẽ rơi xuống mặt đất?
\shortans{3}
\loigiai{Bóng chạm đất khi $h(t)=0\Leftrightarrow -t^2+2t+3=0\Leftrightarrow t=3$ (loại $t=-1$). Vậy sau $3$ giây.}
\end{ex}
